\chapter{Performance Tuning}
\label{ch:performance}

This chapter covers performance optimization for Veilbox — both for the
build process and for runtime operation in virtualized environments.

\section{Build Performance}

\subsection{Parallel Compilation}

The kernel build is the most time-consuming step (approximately 4--5
minutes on a 4-core machine):

\begin{lstlisting}[language=bash]
# Set number of parallel jobs
make -j$(nproc) bzImage
\end{lstlisting}

\subsection{CCache Integration}

CCache can dramatically speed up repeated builds:

\begin{lstlisting}[language=bash]
export CC="ccache gcc"
export CXX="ccache g++"
make -j$(nproc) bzImage
# First build: slow. Subsequent builds: 50-70% faster.
\end{lstlisting}

\subsection{Incremental Builds}

After the first full build, configuration changes rebuild only affected
files:

\begin{lstlisting}[language=bash]
# Changing one kernel config option:
make olddefconfig
make -j$(nproc) bzImage
# Typical time: 30-60 seconds
\end{lstlisting}

\section{Runtime Performance}

\subsection{KVM Acceleration}

QEMU with KVM enables near-native performance:

\begin{lstlisting}[language=bash]
# Check KVM availability
ls -la /dev/kvm
groups  # Must include 'kvm'

# QEMU with KVM
qemu-system-x86_64 -enable-kvm ...
\end{lstlisting}

\subsection{VirtIO Drivers}

VirtIO paravirtualized drivers reduce I/O overhead:

\begin{longtable}{|p{3cm}|p{4cm}|p{5cm}|}
\hline
\textbf{Device} & \textbf{Emulated Alternative} & \textbf{VirtIO Advantage} \\
\hline
Block & IDE (piix3-ide) & 3-5x throughput, 50\% less CPU \\
Network & e1000 / rtl8139 & 5-10x throughput, 30\% less CPU \\
Balloon & None & Dynamic memory management \\
\hline
\end{longtable}

\subsection{Memory Optimization}

\begin{itemize}[nosep]
  \item \textbf{Minimum RAM}: 2 GB. The initramfs occupies ~90 MB,
        containerd uses ~30 MB, kernel uses ~50 MB.
  - \textbf{Transparent Huge Pages}: Enable on host for better memory
        performance.
  - \textbf{Memory overcommit}: Enabled by default for container workloads
        that may overallocate.
\end{itemize}

\section{Kernel Boot Time Optimization}

\subsection{Measured Boot Time}

\begin{lstlisting}
# From QEMU start to login prompt:
# Direct kernel boot: <10 seconds
# GRUB BIOS boot:    <15 seconds
\end{lstlisting}

\subsection{Reducing Boot Time}

\begin{itemize}[nosep]
  \item \textbf{Minimize kernel size}: The stripped-down Veilbox kernel
        is approximately 90 MB (with embedded initramfs). Uncompressed
        parts are proportional to boot time.
  - \textbf{Serial vs. graphical}: Serial console (\texttt{-nographic})
        is faster because no framebuffer initialization is needed.
  - \textbf{Pre-built initramfs}: Skip gen\_initramfs\_list step by
        caching the file list.
\end{itemize}

\section{Container Performance}

\begin{itemize}[nosep]
  \item \textbf{OverlayFS}: Better performance than aufs or device-mapper
        snapshotters. Contains native kernel support.
  - \textbf{noexec /tmp}: Mount /tmp with noexec to prevent container
        escape via writable temp directories.
  - \textbf{ulimit tuning}: Adjust container resource limits with
        \texttt{nerdctl run --ulimit}:
\end{itemize}

\begin{lstlisting}
nerdctl run --ulimit nofile=1024:2048 alpine sh
nerdctl run --memory=256m --cpus=0.5 nginx:alpine
\end{lstlisting}


\section{Kernel Boot Parameters for Performance}

\subsection{mitigations=off}

Disabling CPU vulnerability mitigations can improve performance
significantly in trusted VM environments:

\begin{lstlisting}
# Add to kernel command line:
mitigations=off
\end{lstlisting}

This disables:
\begin{itemize}[nosep]
  \item Spectre v2 (IBPB, IBRS, STIBP)
  \item Meltdown (PTI/KPTI)
  \item Spectre v1 (user pointer sanitization)
  \item Speculative Store Bypass (SSBD)
  \item L1TF (L1D flush)
  \item MDS (microarchitectural data sampling)
\end{itemize}

In a VM where the host is trusted, these mitigations are unnecessary.
Performance gain: 5-15\% on CPU-bound workloads.

\subsection{Scalability Options}

\begin{lstlisting}
# Add to kernel command line:
clocksource=tsc tsc=reliable
\end{lstlisting}

The TSC (Time Stamp Counter) is the fastest available clocksource.
Setting it as the default reduces timer overhead.

\section{I/O Performance}

\subsection{VirtIO-BLK with io\_uring}

The Linux kernel's io\_uring subsystem provides high-performance async
I/O. It can be used by containerd for image pulls:

\begin{lstlisting}
CONFIG_IO_URING=y
\end{lstlisting}

\subsection{Block Device Tuning}

\begin{lstlisting}
# Increase read-ahead buffer
echo 4096 > /sys/block/vda/queue/read_ahead_kb

# Set I/O scheduler (none for VirtIO)
echo none > /sys/block/vda/queue/scheduler

# Increase max sectors per request
echo 2048 > /sys/block/vda/queue/max_sectors_kb
\end{lstlisting}

\section{Memory Performance}

\subsection{Transparent Huge Pages}

THP reduces TLB misses by using 2 MB pages instead of 4 KB:

\begin{lstlisting}
# Enable always (for consistent performance)
echo always > /sys/kernel/mm/transparent_hugepage/enabled

# Defragment when needed
echo defer > /sys/kernel/mm/transparent_hugepage/defrag
\end{lstlisting}

\subsection{Swap}

Swap is not used in Veilbox. The state partition is too small, and
container workloads should fit in RAM:

\begin{lstlisting}
CONFIG_SWAP=n  # Disable swap support in kernel
\end{lstlisting}

\section{Network Performance}

\subsection{Tx Queue Length}

\begin{lstlisting}
# Increase transmit queue length
echo 10000 > /sys/class/net/eth0/tx_queue_len
\end{lstlisting}

\subsection{IRQ Affinity}

Pin network IRQs to specific CPUs:

\begin{lstlisting}
# Find IRQ number for eth0
cat /proc/interrupts | grep eth0

# Set CPU affinity (CPU 0)
echo 1 > /proc/irq/<N>/smp_affinity
\end{lstlisting}
